% Part:sets-functions-relations
% Chapter: size-of-sets
% Section: non-enumerability-alt

\documentclass[../../../include/open-logic-section]{subfiles}

\begin{document}

\olfileid[ar]{sfr}{siz}{nen-alt}

\olsection{مجموعات \printtoken{S}{nonenumerable}}

\begin{editorial}
  يثبت هذا القسم عدم قابلية $\Bin^\omega$ و$\Pow{\Nat}$ للتعداد،
  باستعمال التعريفات الواردة في \olref[enm-alt]{sec}، أي باشتراط
  دالة تقابلية مع~$\Nat$ بدلًا من دالة شاملة من~$\PosInt$.
\end{editorial}

\begin{explain}
مجموعة الأعداد الطبيعية~$\Nat$ لا نهائية، وهي أيضًا، على نحو بديهي،
!!{enumerable}. لكن الحقيقة اللافتة هي أن ثمة مجموعات
\emph{!!{nonenumerable}}، أي مجموعات ليست !!{enumerable}
(انظر \olref[sfr][siz][enm-alt]{defn:enumerable}).

قد يبدو ذلك مفاجئًا. ففي نهاية المطاف، القول إن $A$
!!{nonenumerable} هو القول إنه لا توجد \emph{أي} !!{bijection} $f
\colon \Nat \to A$؛ أي لا توجد دالة ترسل !!{element}s~$\Nat$،
وهي لا نهائية العدد، إلى~$A$ وتستنفد~$A$ كلها. فإذا كانت $A$
!!{nonenumerable}، كان في~$A$ !!{element}s «أكثر» مما في مجموعة
الأعداد الطبيعية.

لإثبات أن مجموعة ما !!{nonenumerable}، لا بد من بيان أنه لا يمكن
وجود !!{bijection} مناسبة. وأفضل سبيل إلى ذلك هو بيان أن كل محاولة
لتعداد !!{element}s~$A$ لا بد أن يفوتها !!{element} واحد على الأقل؛
وهذا يبين أنه لا يمكن لأي دالة $f\colon \Nat \to A$ أن تكون
!!{surjective}. وتتمثل استراتيجية عامة لإثبات ذلك في استعمال
\emph{طريقة كانتور القطرية}. فإذا أعطينا قائمة من !!{element}s~$A$،
ولتكن $x_1$، $x_2$، \dots، أمكننا بناء !!{element} آخر من~$A$ يستحيل،
بحكم طريقة بنائه، أن يكون في تلك القائمة.

لكن أفضل ما يوضح هذا كله هو المثال. فمثالنا الأول هو المجموعة
~$\Bin^\omega$ لجميع السلاسل اللانهائية المؤلفة من~$0$ و~$1$.
(يرمز الرمز «$\Bin$» إلى الثنائية، ويمكننا أن نعده ببساطة المجموعة
ذات العنصرين
$\{0,1\}$.)\oliflabeldef{sfr:card-arithmetic:card-opps:sec}{\footnote{وبعبارة
أدق، ينبغي أن نشترط أن تكون $\Bin^\omega$ مجموعة جميع متتاليات
$\omega$ المؤلفة من~$0$ و~$1$، أي المجموعة
$\funfromto{\omega}{\{0,1\}}$. لكن معنى ذلك لن يتضح إلا في
\olref[sfr][card-arithmetic][card-opps]{sec}.}{} ولا ينبغي، مع ذلك،
أن تسبب هذه الصياغة المتساهلة قليلًا أي التباس في الوقت الحاضر.}
\end{explain}

\begin{thm}
\ollabel{thm:nonenum-bin-omega}
$\Bin^\omega$~!!{nonenumerable}.
\end{thm}

\begin{proof}
تأمل أي تعداد لمجموعة جزئية من~$\Bin^\omega$. فتكون لدينا قائمة ما
$s_{0}$، $s_{1}$، $s_{2}$، \dots{} تكون فيها كل~$s_n$ سلسلة لا
نهائية من~$0$ و~$1$. ولتكن $s_n(m)$ الرقم ذا الرتبة~$n$ من السلسلة
ذات الرتبة~$m$ في هذه القائمة. ويمكننا إذن أن نتصور قائمتنا مصفوفةً،
حيث يوضع $s_n(m)$ في الصف ذي الرتبة~$n$ والعمود ذي الرتبة~$m$:
\[
\begin{array}{c|c|c|c|c|c}
& 0 & 1 & 2 & 3 & \dots \\\hline
0 & \mathbf{s_{0}(0)} & s_{0}(1) & s_{0}(2) & s_0(3) & \dots \\\hline
1 & s_{1}(0)& \mathbf{s_{1}(1)} & s_1(2) & s_1(3) & \dots \\\hline
2 & s_{2}(0)& s_{2}(1) & \mathbf{s_2(2)} & s_2(3) & \dots \\\hline
3 & s_{3}(0)& s_{3}(1) & s_3(2) & \mathbf{s_3(3)} & \dots \\\hline
\vdots & \vdots & \vdots & \vdots & \vdots & \mathbf{\ddots}
\end{array}
\]
وسنبني الآن سلسلة لا نهائية~$d$ من~$0$ و~$1$ ليست في هذه القائمة.
ونفعل ذلك بتعيين كل خانة من خاناتها، أي بتعيين $d(n)$ لكل $n \in
\Nat$. ومن الوجهة الحدسية، نقرأ قطر المصفوفة أعلاه من أوله إلى آخره
(ومن هنا اسم «الطريقة القطرية»)، ثم نغير كل~$1$ إلى~$0$ وكل~$1$
إلى~$0$. وبصورة أكثر تجريدًا، نعرّف $d(n)$ ليكون $0$ أو~$1$ بحسب ما
إذا كان !!{element} ذو الرتبة~$n$ على القطر، أي $s_n(n)$، هو $1$
أو~$0$، أي:
\[
d(n) =
\begin{cases}
1 & \text{إذا كان $s_{n}(n) = 0$}\\
0 & \text{إذا كان $s_{n}(n) = 1$}
\end{cases}
\]
ومن الواضح أن $d \in \Bin^\omega$، لأنها سلسلة لا نهائية من~$0$
و~$1$. لكننا بنينا~$d$ بحيث يكون $d(n) \neq s_n(n)$ لكل $n \in
\Nat$. أي إن~$d$ تختلف عن~$s_n$ في خانتها ذات الرتبة~$n$.
ومن ثم يكون $d \neq s_n$ لكل $n\in \Nat$. ولذلك لا يمكن أن تكون~$d$
في القائمة $s_0$، $s_1$، $s_2$،
\dots

لقد بينا، إذا أعطينا تعدادًا اعتباطيًا لمجموعة جزئية ما من
$\Bin^\omega$، أن !!{element} من~$\Bin^\omega$ سيفوت ذلك التعداد. ولذلك لا
يوجد تعداد للمجموعة~$\Bin^\omega$، أي إن~$\Bin^\omega$
!!{nonenumerable}.
\end{proof}

\begin{explain}
تسمى طريقة البرهان هذه «الاستدلال القطري» لأنها تستعمل قطر المصفوفة
لتعريف~$d$. لكن الاستدلال القطري لا يستلزم وجود مصفوفة. إذ يمكننا،
في الواقع، أن نبين أن مجموعة ما !!{nonenumerable} باستعمال فكرة
مماثلة، حتى حين لا تكون ثمة مصفوفة ولا قطر فعلي. وتوضح النتيجة الآتية
كيفية ذلك.
\end{explain}

\begin{thm}
\ollabel{thm:nonenum-pownat}
$\Pow{\Nat}$ ليست !!{enumerable}.
\end{thm}

\begin{proof}
نسير بالطريقة نفسها، فنبين أن كل قائمة من المجموعات الجزئية من~$\Nat$
تغفل مجموعة جزئية ما من~$\Nat$. افترض إذن أن لدينا قائمة ما
$N_0, N_1, N_2, \ldots$ من المجموعات الجزئية من~$\Nat$. ونعرّف
مجموعة~$D$ على النحو الآتي: $n \in D$ إذا وفقط إذا كان $n \notin
N_{n}$:
\[
D = \Setabs{n \in \Nat}{n \notin N_n}
\]
من الواضح أن $D\subseteq \Nat$. لكن لا يمكن أن تكون~$D$ في القائمة.
فبحكم بنائها، $n \in D$ إذا وفقط إذا كان $n\notin N_n$، ولذلك يكون
$D \neq N_n$ لكل $n \in \Nat$.
\end{proof}

\begin{explain}
لم يذكر البرهان السابق قطرًا، لكن يمكنك أن تتصوره متضمنًا قطرًا إذا
صورته على النحو الآتي: تصور المجموعات $N_0$، $N_1$، \dots، مكتوبةً
في مصفوفة، بحيث نكتب $N_n$ في الصف ذي الرتبة~$n$ بكتابة~$m$ في
العمود ذي الرتبة~$m$ إذا وفقط إذا كان $m \in N_n$. ولتكن، مثلًا،
المجموعات الأربع الأولى في تلك القائمة
$\{0,1,2,\dots\}$، $\{1, 3, 5, \dots\}$، $\{0,1,4\}$، و
$\{2,3,4,\dots\}$؛ فعندئذ تبدأ المصفوفة هكذا:
\[
\begin{array}{r@{}rrrrrrr}
  N_0 = \{ & \mathbf{0}, & 1, & 2, & & & & \dots\}\\
  N_1 = \{ &  & \mathbf{1}, &  & 3, &  & 5, & \dots\}\\
  N_2 = \{ & 0, & 1, &  &  & 4\phantom{,} &  & \}\\
  N_3 = \{ &  &  & 2, & \mathbf{3}, & 4, & & \dots\}\\
  &\vdots & & & & & & \ddots\phantom{\}}
  \end{array}
\]
وعندئذ تكون~$D$ المجموعة الناتجة من النزول على القطر، بحيث نضع
$n \in D$ إذا وفقط إذا كان~$n$ \emph{غير موجود} على القطر. ففي
الحالة أعلاه نحذف~$0$ و~$1$، وندرج~$2$، ونحذف~$3$، وهكذا.
\end{explain}

\begin{prob}
أثبت، بحجة قطرية صريحة، أن مجموعة جميع الدوال $f \colon \Nat \to
\Nat$ !!{nonenumerable}. أي أثبت أنه إذا كانت $f_1$، $f_2$، \dots،
قائمة من الدوال، وكانت كل $f_i\colon \Nat \to \Nat$، فثمة دالة
$g \colon \Nat \to \Nat$ ليست في هذه القائمة.
\end{prob}

\end{document}
